\begin{textsample}{NEG Dim 1 – 2017\_c4 – Score -4.00 – 2017\_c4/enade\_2017\_qui\_5a144fb539.txt}  \label{ex:f1_neg_168}
CAPÍTULO 4 PERCEPÇÃO DA PROVA As análises feitas neste capítulo tratam das percepções dos concluintes da Área de Química sobre a prova aplicada no Enade/2017.
Estas percepções foram mensuradas por meio de nove questões que avaliaram desde o grau de dificuldade da prova até o tempo gasto para concluí -la.
As percepções sobre a prova foram relacionadas ao desempenho dos estudantes e à Grande Região de funcionamento do curso.
O questionário de percepção da prova encontra -se ao final do Anexo VII, que traz a reprodução da prova.
O desempenho dos estudantes foi classificado em quatro quartos.
Para tanto, esse desempenho foi ordenado de forma ascendente.
O percentil 25, P25, também conhecido como primeiro quartil, é a nota de desempenho que deixa um quarto ( 25 \% ) dos valores observados abaixo e três quartos acima.
A Figura 1 apresenta uma ilustração deste conceito.
O quarto inferior de desempenho é composto pelas notas abaixo do primeiro quartil.
Já o percentil 75, P75, também conhecido como terceiro quartil, é o valor para o qual há três quartos ( 75 \% ) dos dados abaixo e um quarto acima dele.
O quarto superior de desempenho é composto pelas notas iguais ou acima do terceiro quartil.
O percentil 50, P50, também conhecido como mediana, é o valor que divide as notas em dois conjuntos de igual tamanho.
O segundo quarto inclui valores entre o primeiro quartil ( P25 ) e a mediana.
O terceiro quarto contém os valores entre a mediana ( P50 ) e o terceiro quartil ( P75 ).
Vale ressaltar que percentis, quartis e medianas são pontos que não obrigatoriamente pertencem ao conjunto original de dados, ao passo que os quartos são subconjuntos dos dados originais.
A seguir, serão apresentados gráficos com resultados selecionados, relativos às nove questões avaliadas por grupos de estudantes.
Os gráficos apresentam nas barras o percentual de alunos que assinalaram uma das opções ou a soma das porcentagens daqueles que assinalaram duas ( ou três ) delas.
Por exemplo, para as questões 1 e 2, os gráficos apresentam a porcentagem total de participantes que assinalaram as opções ( D ) Difícil e ( E ) Muito difícil.
Em \textbf{cada} barra foram assinalados também os extremos do intervalo de confiança de 95 \% como linhas verticais unidas por uma linha horizontal na forma da letra H maiúscula.
As Tabelas no Anexo II apresentam os valores absolutos e a distribuição percentual17 das alternativas válidas das nove questões, segundo o mesmo recorte de desempenho dos alunos e Grande Região de funcionamento do curso. 4.1 GRAU DE DIFICULDADE DA PROVA 4.1.1 Componente de Formação Geral Ao avaliarem Qual o grau de dificuldade desta prova na parte de Formação Geral? ( Questão 1 ), 26,9 \% do grupo de inscritos e presentes de Licenciatura optaram pelas alternativas Difícil ou Muito difícil.
Entretanto, para mais da metade dos estudantes ( 62,6 \% ), o Componente de Formação Geral da prova foi considerado com grau de dificuldade Médio ( Gráfico 4.1a, Gráfico 4.2a e, no Anexo II, a Tabela II.1a ).
O percentual de estudantes de Licenciatura que consideraram a prova como Difícil ou Muito difícil foi maior na região Norte, onde a proporção foi de 33,7 \%, enquanto a de menor incidência foi a Sudeste, com 23,8 \%.
No Gráfico 4.1a, é possível observar que as diferenças das regiões Nordeste e a Sudeste para as regiões Norte e Centro-Oeste são estatisticamente significativas.
Nas Grandes Regiões, a proporção de presentes à prova que consideraram o Componente de Formação Geral como sendo de grau de dificuldade Médio \textbf{esteve} entre 57,1 \%, na região Norte, e 64,1 \%, na região Sul. 17 Cumpre \textbf{lembrar} uma das convenções para tabelas numéricas ( pág. iii ) sobre a possibilidade de a soma das partes não resultar em 100 \% por questões de arredondamento.
Ao avaliarem Qual o grau de dificuldade desta prova na parte de Formação Geral? ( Questão 1 ), 20,0 \% do grupo de inscritos e presentes de Bacharelado optaram pelas alternativas Difícil ou Muito difícil.
Entretanto, para mais da metade dos estudantes ( 62,3 \% ), o Componente de Formação Geral da prova foi considerado com grau de dificuldade Médio ( Gráfico 4.1b, Gráfico 4.2b e, no Anexo II, a Tabela II.1b ).
O percentual de estudantes de Bacharelado que consideraram a prova Difícil ou Muito difícil foi maior na região Norte, onde a proporção foi de 33,3 \%, enquanto a de menor incidência foi a Nordeste, com 18,9 \%.
No Gráfico 4.1b, é possível observar que as diferenças entre as Grandes regiões não são estatisticamente significativas.
Nas Grandes Regiões, a proporção de presentes à prova que consideraram o Componente de Formação Geral como sendo de grau de dificuldade Médio \textbf{esteve} entre 51,9 \%, na região Norte, e 64,2 \% na região Nordeste.
O percentual de alunos de Licenciatura que consideraram a prova Difícil ou Muito difícil foi decrescente em função dos quartos de desempenho : 34,0 \% no primeiro quarto e 22,2 \% no quarto, grupo de melhor desempenho na prova.
Nos quartos de desempenho intermediários, a proporção de alunos que consideraram a prova Difícil ou Muito difícil foi de 29,1 \% no segundo quarto e 22,5 \% no terceiro quarto.
As diferenças entre qualquer um dos dois primeiros quartos e os dois últimos são estatisticamente significativas.
Para todos os quartos de desempenho, a alternativa modal para esta pergunta foi Médio, com 58,7 \% e 64,0 \% dos respondentes nos quartos extremos, primeiro e quarto, respectivamente.
O percentual de alunos de Bacharelado que consideraram a prova Difícil ou Muito difícil foi decrescente em função dos quartos de desempenho : 27,5 \% no primeiro quarto e 14,2 \% no quarto, grupo de melhor desempenho na prova.
Nos quartos de desempenho intermediários, a proporção de alunos que consideraram a prova Difícil ou Muito difícil foi de 23,9 \% no segundo quarto e 14,7 \% no terceiro quarto.
As diferenças entre as proporções dos dois primeiros quartos a as dos dois últimos são estatisticamente significativas.
Para todos os quartos de desempenho, a alternativa modal para esta pergunta foi Médio, com 59,9 \% e 63,0 \% dos respondentes nos quartos extremos, primeiro e quarto, respectivamente. 4.1.2 Componente de Conhecimento Específico Ao responderem à Questão 2 – Qual o grau de dificuldade desta prova na parte de Componente Específico? – 54,9 \% do grupo de estudantes de Licenciatura classificaram -na como Difícil ou Muito difícil.
Além disso, o Componente de Conhecimento Específico da prova foi considerado com grau de dificuldade Médio por 42,5 \% dos alunos ( Gráfico 4.3a, Gráfico 4.4a, e, no Anexo II, a Tabela II.2a ).
A análise das respostas dos estudantes de Licenciatura quanto ao grau de dificuldade do Componente de Conhecimento Específico da prova, agregado por Grande Região, mostra que as diferenças entre a maior e as duas menores proporções de alunos que a avaliaram como Difícil ou Muito difícil são estatisticamente significativas : a maior na região Centro-Oeste ( 61,0 \% ) e as duas menores proporções nas regiões Sudeste ( 51,5 \% ) e Sul ( 50,9 \% ).
O percentual de alunos que classificaram o grau de dificuldade como Médio, no Componente de Conhecimento Específico, variou de 36,0 \% a 46,1 \%, para as regiões Norte e Sudeste, respectivamente.
Ao responderem à Questão 2 – Qual o grau de dificuldade desta prova na parte de Componente Específico? – 55,7 \% do grupo de estudantes de Bacharelado classificaram -na como Difícil ou Muito difícil.
Além disso, o Componente de Conhecimento Específico da prova foi considerado com grau de dificuldade Médio por 41,6 \% dos alunos ( Gráfico 4.3b, Gráfico 4.4b, e, no Anexo II, a Tabela II.2b ).
A análise das respostas dos estudantes de Bacharelado quanto ao grau de dificuldade do Componente de Conhecimento Específico da prova, agregado por Grande Região, mostra que não há diferença estatisticamente significativa entre as Grandes Regiões.
O percentual de alunos que classificaram o grau de dificuldade como Médio, no Componente de Conhecimento Específico, variou de 38,7 \% a 47,5 \%, para as regiões Centro-Oeste e Nordeste, respectivamente.
Considerando -se a avaliação da dificuldade das questões do Componente de Conhecimento Específico da prova, de acordo com o desempenho dos estudantes de Licenciatura, observa -se que não há diferença estatisticamente significativa entre os resultados das Grandes Regiões.
A proporção dos que classificaram a parte específica como Difícil ou Muito difícil variou de 52,5 \% ( último quarto ) a 57,8 \% ( segundo quarto ).
A alternativa com a segunda maior proporção de escolha em todos os quartos para a Questão 2 foi o grau Médio, com 42,9 \% do quarto inferior e 44,5 \% do superior optando por esta resposta.
Considerando -se a avaliação da dificuldade das questões do Componente de Conhecimento Específico da prova, de acordo com o desempenho dos estudantes de Bacharelado, observa -se que há diferença estatisticamente significativa dos resultados entre o último e os demais quartos de desempenho.
A proporção dos que classificaram a parte específica como Difícil ou Muito difícil variou de 46,4 \% ( último quarto ) a 61,8 \% ( segundo quarto ).
A alternativa modal no último quarto para a Questão 2 foi o grau Médio, sendo essa a segunda alternativa com maior proporção nos outros três quartos.
Tal alternativa variou de 35,8 \% no segundo quarto a 48,9 \% no superior. 4.2 EXTENSÃO DA PROVA EM RELAÇÃO AO TEMPO TOTAL Indagados quanto à extensão da prova, em relação ao tempo total oferecido para a sua resolução ( Questão 3 ), os estudantes de Licenciatura apontaram, com maior incidência, a alternativa que considerava a extensão adequada, para todas as agregações consideradas ( Gráfico 4.5a, Gráfico 4.6a, e, no Anexo II, a Tabela II.3a ).
O percentual de alunos de Licenciatura que responderam ser a extensão da prova adequada foi de 57,7 \%.
Já 36,1 \% dos inscritos presentes consideraram que a prova foi longa ou muito longa, e 6,2 \% a avaliaram como curta ou muito curta.
Entre as Grandes Regiões, a proporção daqueles que avaliaram a prova como longa ou muito longa, em relação ao tempo total destinado à sua resolução variou de 25,1 \% na Entre as Grandes Regiões, a proporção daqueles que avaliaram a prova como longa ou muito longa, em relação ao tempo total destinado à sua resolução variou de 25,1 \% na região Norte até 38,2 \% na região Sudeste.
As diferenças entre a região Norte e as demais são estatisticamente significativas.
Indagados quanto à extensão da prova, em relação ao tempo total oferecido para a sua resolução ( Questão 3 ), os estudantes de Bacharelado apontaram, com maior incidência, a alternativa que considerava a extensão adequada, para todas as agregações consideradas ( Gráfico 4.5b, Gráfico 4.6b, e, no Anexo II, a Tabela II.3b ).
O percentual de alunos de Bacharelado que responderam ser a extensão da prova adequada foi de 58,0 \%.
Já 36,9 \% dos inscritos presentes consideraram que a prova foi longa ou muito longa, e 5,0 \% a avaliaram como curta ou muito curta.
Entre as Grandes Regiões, a proporção daqueles que avaliaram a prova como longa ou muito longa, em relação ao tempo total destinado à sua resolução, variou de 29,6 \% na região Norte até 48,1 \% na região Nordeste.
Só há diferença estatisticamente significativa entre a região Nordeste e a Sudeste.
Considerando -se o desempenho dos alunos de Licenciatura, nota -se ainda que 58,2 \% consideraram a extensão da prova adequada no quarto de desempenho inferior e 56,7 \% no de melhor desempenho ( quarto superior ).
Nos quartos intermediários, esta proporção foi 57,0 \% no segundo quarto e 59,1 \% no terceiro.
No Gráfico 4.6a, pode -se constatar que, grosso modo, há uma tendência crescente da proporção de estudantes de Licenciatura que consideraram a prova longa ou muito longa em função dos quartos de desempenho.
De fato, esse crescimento não ocorre do segundo para o terceiro quarto, mas, nas comparações dos quartos de desempenho, o crescimento é observado.
Entretanto, não há diferença estatisticamente significativa entre as proporções de escolha desta resposta.
Considerando -se o desempenho dos alunos de Bacharelado, nota -se ainda que 34,2 \% consideraram a extensão da prova adequada no quarto de desempenho inferior e 39,6 \% no de melhor desempenho ( quarto superior ).
Nos quartos intermediários, esta proporção foi 34,8 \% no segundo quarto e 39,1 \% no terceiro.
No Gráfico 4.6b, pode -se constatar que há uma tendência crescente da proporção de estudantes de Bacharelado que consideraram a prova longa ou muito longa em função dos quartos de desempenho.
Observa -se, porém, que não há diferença estatisticamente significativa entre as proporções de escolha desta resposta. 4.3 COMPREENSÃO DOS ENUNCIADOS DAS QUESTÕES 4.3.1 Componente de Formação Geral Com relação aos enunciados das questões do Componente de Formação Geral ( Questão 4 ), as opiniões foram positivas, já que 78,2 \% dos alunos de Licenciatura avaliados consideraram todos ou a maioria dos enunciados das questões claros e objetivos ( Gráfico 4.7a, Gráfico 4.8a, e, no Anexo II, a Tabela II.4a ).
Na análise regional, a percentagem de estudantes de Licenciatura que avaliaram que todos ou a maioria dos enunciados das questões do Componente de Formação Geral \textbf{estavam} claros e objetivos variou de 75,0 \% na região Nordeste a 81,7 \% na região Centro-Oeste, sendo esta diferença estatisticamente significativa.
A análise das percepções dos estudantes de Licenciatura sobre a clareza e objetividade dos enunciados permite afirmar que todos, ou a maioria dos enunciados de questões relativas ao Componente de Formação Geral, foram considerados claros e objetivos para a maior parte dos respondentes ( maior ou igual a 75,0 \% em todas as regiões e maior ou igual a 70,5 \% para todos os quartos de desempenho ).
Com relação aos enunciados das questões do Componente de Formação Geral ( Questão 4 ), as opiniões foram positivas, já que 80,0 \% dos alunos de Bacharelado avaliados consideraram todos ou a maioria dos enunciados das questões claros e objetivos ( Gráfico 4.7b, Gráfico 4.8b, e, no Anexo II, a Tabela II.4b ).
Na análise regional, a percentagem de estudantes de Bacharelado que avaliaram que todos ou a maioria dos enunciados das questões do Componente de Formação Geral \textbf{estavam} claros e objetivos variou de 78,2 \% na região Nordeste a 83,8 \% na região Centro-Oeste, não sendo, entretanto, esta diferença estatisticamente significativa.
A análise das percepções dos estudantes de Bacharelado sobre a clareza e objetividade dos enunciados permite afirmar que todos, ou a maioria dos enunciados de questões relativas ao Componente de Formação Geral, foram considerados claros e objetivos para a maior parte dos respondentes ( maior ou igual a 78,2 \% em todas as regiões e maior ou igual a 76,8 \% para todos os quartos de desempenho ).
Segundo o desempenho, observa -se que a proporção dos alunos de Licenciatura que emitiram esta opinião cresce conforme o desempenho aumenta, com diferença estatisticamente significativas entre o primeiro e o segundo quartos de desempenho.
No quarto superior, a clareza e objetividade de todos ou da maioria dos enunciados das questões foi percebida por 84,5 \%.
Segundo o desempenho, observa -se que a proporção dos alunos de Bacharelado que emitiram esta opinião cresce conforme o desempenho aumenta nos três primeiros quartos de desempenho, sem, entretanto, diferenças estatisticamente significativas.
No terceiro quarto de desempenho, a clareza e objetividade de todos ou da maioria dos enunciados das questões foi percebida por 82,2 \% dos participantes, e no quarto superior, o percentual foi de 81,0 \%. 4.3.2 Componente de Conhecimento Específico Com relação aos enunciados das questões do Componente de Conhecimento Específico da prova, para 76,2 \% dos estudantes de Licenciatura avaliados da Área de Química, a clareza e a objetividade ( Questão 5 ) \textbf{estavam} presentes em todas ou na maioria das questões ( Gráfico 4.9a, Gráfico 4.10a, e no Anexo II, a Tabela II.5a ).
A maioria dos estudantes de Licenciatura de todas as Grandes Regiões brasileiras considerou claros e objetivos todos ou a maioria dos enunciados das questões do Componente de Conhecimento Específico da prova, percentual sempre maior ou igual a 71,3 \%.
As diferenças entre algumas das Grandes Regiões são estatisticamente significativas.
Com relação aos enunciados das questões do Componente de Conhecimento Específico da prova, para 79,4 \% dos estudantes de Bacharelado avaliados da Área de Química, a clareza e a objetividade ( Questão 5 ) \textbf{estavam} presentes em todas ou na maioria das questões ( Gráfico 4.9b, Gráfico 4.10b, e no Anexo II, a Tabela II.5b ).
A maioria dos estudantes de Bacharelado de todas as Grandes Regiões brasileiras considerou claros e objetivos todos ou a maioria dos enunciados das questões do Componente de Conhecimento Específico da prova, percentual sempre maior ou igual a 73,7 \%.
Não há diferença estatisticamente significativa entre as Grandes Regiões.
A proporção de estudantes de Licenciatura que consideraram os enunciados das questões claros e objetivos apresenta uma tendência crescente em relação ao aumento de desempenho : mais elevada no quarto superior ( 83,0 \% ) se comparada ao quarto inferior de desempenho ( 66,8 \% ).
As diferenças entre o primeiro quarto, o de menor frequência, o os demais quartos são estatisticamente significativas.
A proporção de estudantes de Bacharelado que consideraram os enunciados das questões claros e objetivos apresenta uma tendência crescente em relação ao aumento de desempenho : mais elevada no quarto superior ( 84,8 \% ) se comparada ao quarto inferior de desempenho ( 72,0 \% ).
As diferenças entre o primeiro quarto, o de menor frequência, e os demais quartos são estatisticamente significativas. 4.4 SUFICIÊNCIA DAS INFORMAÇÕES/INSTRUÇÕES FORNECIDAS Ao avaliarem as informações/instruções fornecidas para a resolução das questões ( Questão 6 ), 79,5 \% dos respondentes de Licenciatura da Área de Química de todo o Brasil afirmaram que estas eram até excessivas ou suficientes em todas ou na maioria das questões ( Gráfico 4.11a, Gráfico 4.12a, e, no Anexo II, a Tabela II.6a ).
Quanto à distribuição de respondentes pelas Grandes Regiões observa -se que a proporção de estudantes que consideraram as informações/instruções fornecidas até excessivas ou suficientes em todas ou na maioria das questões foi sempre superior ou igual a 72,1 \% ( região Norte ) chegando a 84,3 \% na região Sul.
A diferença entre o maior percentual, obtido na região Norte, e o menor, da região Sul, é estatisticamente significativa.
Ao avaliarem as informações/instruções fornecidas para a resolução das questões ( Questão 6 ), 81,8 \% dos respondentes de Bacharelado da Área de Química de todo o Brasil afirmaram que estas eram até excessivas ou suficientes em todas ou na maioria das questões ( Gráfico 4.11b, Gráfico 4.12b, e, no Anexo II, a Tabela II.6b ).
Quanto à distribuição de respondentes pelas Grandes Regiões, observa -se que a proporção de estudantes que consideraram as informações/instruções fornecidas até excessivas ou suficientes, em todas ou na maioria das questões, foi sempre superior ou igual a 72,2 \%, chegando a 82,7 \% na região Sul.
Não há diferenças estatisticamente significativas entre as Grandes Regiões.
Levando -se em conta o desempenho dos participantes de Licenciatura, notam -se diferenças estatisticamente significativas das opiniões de estudantes de Licenciatura de entre os três últimos quartos de desempenho, e entre o primeiro e os dois últimos, como mostra o Gráfico 4.12a.
O percentual de participantes que avaliaram as informações/instruções como até excessivas ou suficientes, em todas ou na maioria das questões, foi mais elevado no quarto superior ( 86,5 \% ), percentual superior à média nacional ( 79,5 \% ).
Já no quarto inferior, a suficiência das informações/instruções declarada como até excessiva, em todas ou na maioria das questões, foi percebida por 74,7 \% dos respondentes.
Estas proporções são crescentes com os quartos de desempenho, menos do primeiro para o segundo quarto, cujas proporções são 74,7 \% e 74,6 \%, respectivamente.
Levando -se em conta o desempenho dos participantes de Bacharelado, nota -se diferença estatisticamente significativa entre as opiniões de estudantes de Bacharelado dos quartos extremos de desempenho, primeiro e último, como mostra o Gráfico 4.12b.
O percentual de participantes que avaliaram as informações/instruções como até excessivas ou suficientes, em todas ou na maioria das questões, foi mais elevado no quarto superior ( 88,0 \% ), percentual superior à média nacional ( 81,8 \% ).
Já no quarto inferior, a suficiência das informações/instruções declaradas como até excessiva, em todas ou na maioria das questões foi percebida por 78,3 \% dos respondentes.
Estas proporções são crescentes com os quartos de desempenho, menos do primeiro para o segundo quarto, cujas proporções são 78,3 \% e 77,9 \%, respectivamente. 4.5 DIFICULDADE ENCONTRADA AO RESPONDER À PROVA \textbf{Perguntados} sobre as dificuldades com as quais depararam ao responder à prova ( Questão 7 ), 21,8 \% dos estudantes de Licenciatura apontaram o Desconhecimento do conteúdo.
Para 45,6 \%, a Forma diferente de abordagem do conteúdo foi indicada como dificuldade.
Já a Falta de motivação para fazer a prova foi a dificuldade apontada por 18,9 \% dos respondentes.
Considerando -se todo o Brasil, 8,9 \% dos respondentes afirmaram que não tiveram qualquer tipo de dificuldade para responder à prova ( Tabela II.7a no Anexo II ).
Os Gráficos 4.13a e 4.14a apresentam os percentuais de estudantes de Licenciatura que apontaram o Desconhecimento do conteúdo como dificuldade percebida ao responder à prova.
Na análise por Grandes Regiões, o percentual de inscritos e presentes de Licenciatura que apontaram o Desconhecimento do conteúdo como dificuldade ao responder à prova não superou 24,4 \%.
Os percentuais variaram de 20,1 \% na região Nordeste a 24,4 \% na Centro-Oeste, não sendo, entretanto, esta diferença significativa estatisticamente.
A Forma diferente de abordagem do conteúdo foi a escolha modal dos estudantes de Licenciatura em todas as Grandes Regiões, com percentuais que variaram de 41,2 \% ( região Sudeste ) a 51,5 \% ( Sul ).
O percentual de alunos que citaram a Falta de motivação como dificuldade variou de 14,7 \% ( região Centro-Oeste ) a 19,9 \% ( regiões Nordeste e Sudeste ).
Os que declararam não ter qualquer tipo de dificuldade para responder à prova variaram de 6,5 \% na região Nordeste a 11,4 \% na Sudeste. \textbf{Perguntados} sobre as dificuldades com as quais depararam ao responder à prova ( Questão 7 ), 22,2 \% dos estudantes de Bacharelado apontaram o Desconhecimento do conteúdo.
Para 44,9 \%, a Forma diferente de abordagem do conteúdo foi indicada como dificuldade.
Já a Falta de motivação para fazer a prova foi a dificuldade apontada por 18,9 \% dos respondentes.
Considerando -se todo o Brasil, 10,4 \% dos respondentes afirmaram que não tiveram qualquer tipo de dificuldade para responder à prova ( Tabela II.7b no Anexo II ).
Os Gráficos 4.13b e 4.14b apresentam os percentuais de estudantes de Bacharelado que apontaram o Desconhecimento do conteúdo como dificuldade percebida ao responder à prova.
Na análise por Grandes Regiões, o percentual de inscritos e presentes de Bacharelado que apontaram o Desconhecimento do conteúdo como dificuldade ao responder à prova não superou 28,6 \%.
Os percentuais variaram de 18,9 \% na região Sul a 28,6 \% na Centro-Oeste, sendo esta diferença significativa estatisticamente.
A Forma diferente de abordagem do conteúdo foi a escolha modal dos estudantes de Bacharelado em todas as regiões, com percentuais que variaram de 31,5 \% ( região Norte ) a 49,4 \% ( região Sul ).
O percentual de alunos que citaram a Falta de motivação como dificuldade variou de 16,7 \% ( região Norte ) a 24,7 \% ( região Nordeste ).
Os que declararam não ter qualquer tipo de dificuldade para responder à prova variaram de 6,2 \% na região Nordeste a 16,7 \% na Norte.
Com relação aos quartos de desempenho, o Desconhecimento do conteúdo foi a opção escolhida por 25,0 \% dos estudantes de Licenciatura do quarto superior e 18,4 \% do quarto inferior.
Nota -se que a proporção do primeiro quarto apresenta diferença estatisticamente significativa para o último.
A alternativa modal para os alunos de Licenciatura, quando agregados pelos quartos de desempenho, foi que a dificuldade encontrada foi causada pela Forma diferente de abordagem do conteúdo : 47,1 \% no quarto inferior e 41,0 \% do quarto superior assim o responderam.
Com relação aos quartos de desempenho, o Desconhecimento do conteúdo foi a opção escolhida por 27,5 \% dos estudantes de Bacharelado do quarto superior e 18,0 \% do quarto inferior.
Nota -se que a proporção do primeiro quarto apresenta diferença estatisticamente significativa para o último.
A alternativa modal para os alunos de Bacharelado, quando agregados pelos quartos de desempenho, foi que a dificuldade encontrada foi causada pela Forma diferente de abordagem do conteúdo : 47,5 \% no quarto inferior e 37,7 \% do quarto superior assim o responderam. 4.6 CONTEÚDOS DAS QUESTÕES OBJETIVAS DA PROVA Ao analisar os conteúdos das questões objetivas da prova ( Questão 8 ), um percentual muito pequeno dos estudantes de Licenciaturas avaliados, apenas 5,4 \%, afirmou que não estudou ainda a maioria desses conteúdos ( Gráficos 4.15a, Gráfico 4.16a, e a Tabelas II.8a no Anexo II ).
A maioria ( 54,4 \% ) afirmou ter estudado e aprendido muitos ou todos os conteúdos avaliados.
Na análise por Grande Região, a proporção de respondentes de Licenciatura que escolheram a opção não estudou ainda a maioria desses conteúdos foi pequena.
Observa -se que nas regiões Norte ( 9,6 \% ), Nordeste ( 6,1 \% ) e Centro-Oeste ( 5,6 \% ), apesar de pequenas, as proporções foram maiores do que a média nacional ( 5,4 \% ).
Observa -se diferença estatisticamente significativa entre algumas das regiões.
Em todas as regiões, a maioria dos presentes afirmou ter estudado e aprendido muitos ou todos os conteúdos, com proporções variando entre 41,3 \% na região Norte e 64,2 \% na Sudeste.
Ao analisar os conteúdos das questões objetivas da prova ( Questão 8 ), um percentual muito pequeno dos estudantes de Bacharelados avaliados, apenas 3,5 \%, afirmou que não estudou ainda a maioria desses conteúdos ( Gráficos 4.15b, Gráfico 4.16b, e a Tabelas II.8b no Anexo II ).
A maioria ( 63,9 \% ) afirmou ter estudado e aprendido muitos ou todos os conteúdos avaliados.
Na análise por Grande Região, a proporção de respondentes de Bacharelado que escolheram a opção não estudou ainda a maioria desses conteúdos foi pequena.
Observa -se que nas regiões Norte ( 7,5 \% ), Nordeste ( 4,5 \% ) e Centro-Oeste ( 4,4 \% ), apesar de pequenas, as proporções foram maiores do que a média nacional ( 3,5 \% ).
Não observa -se diferença estatisticamente significativa entre as regiões.
Em todas as regiões, a maioria dos presentes afirmou ter estudado e aprendido muitos ou todos os conteúdos, com proporções variando entre 52,6 \% na região Centro-Oeste e 67,2 \% na Sul.
Considerando -se separadamente as opiniões de estudantes de Licenciatura dos quatro quartos de desempenho, observa -se que, no quarto inferior, 9,6 \% ofereceram como resposta não estudou ainda a maioria desses conteúdos, sendo 1,9 \% os do quarto superior com a mesma resposta.
A diferença entre os alunos que optaram por este motivo de dificuldade entre os quartos extremos de desempenho é estatisticamente significativa, a exceção da diferença entre o primeiro e o segundo quarto.
Tendo em conta o quarto superior, 69,4 \% dos alunos de Licenciatura afirmaram ter estudado e aprendido muitos ou todos os conteúdos.
No outro extremo, no primeiro quarto, 40,2 \% optaram pelas mesmas categorias.
As proporções são crescentes com o desempenho.
Considerando -se separadamente as opiniões de estudantes de Bacharelado dos quatro quartos de desempenho, observa -se que, no quarto inferior, 4,4 \% ofereceram como resposta que não estudou ainda a maioria desses conteúdos, sendo 2,1 \% os do quarto superior com a mesma resposta.
As diferenças entre os alunos que optaram por este motivo de dificuldade entre os quartos de desempenho não são estatisticamente significativas.
Tendo em conta o quarto superior, 76,1 \% dos alunos de Bacharelado afirmaram ter estudado e aprendido muitos ou todos os conteúdos.
No outro extremo, no primeiro quarto, 54,7 \% optaram pelas mesmas categorias.
As proporções são crescentes com o desempenho. 4.7 TEMPO GASTO PARA CONCLUIR A PROVA Ao responderem sobre o tempo de conclusão da prova ( Questão 9 ), mais do que três quartos dos estudantes de Licenciatura ( 77,2 \% ) afirmaram ter gasto Entre duas e quatro horas ( Gráfico 4.17a, Gráfico 4.18a e, no Anexo II, a Tabela II.9a ).
Considerando -se as cinco Grandes Regiões brasileiras, os percentuais dos que utilizaram Entre duas e quatro horas para finalizar a prova nas regiões Norte ( 72,0 \% ) e Nordeste ( 75,4 \% ) foram inferiores ao percentual nacional.
Nas demais Grandes Regiões, o percentual de alunos de Licenciatura que dispensaram Entre duas e quatro horas para concluir a prova ficou igual ou acima de 77,6 \%, como mostra o Gráfico 4.17a.
Há diferença estatisticamente significativa entre a proporção do Norte e a da região Sul, a menor e a maior proporção, respectivamente.
Ao responderem sobre o tempo de conclusão da prova ( Questão 9 ), mais do que três quartos dos estudantes de Bacharelado ( 77,6 \% ) afirmaram ter gasto Entre duas e quatro horas ( Gráfico 4.17b, Gráfico 4.18b e, no Anexo II, a Tabela II.9b ).
Considerando -se as cinco Grandes Regiões brasileiras, os percentuais dos que utilizaram Entre duas e quatro horas para finalizar a prova nas regiões Norte ( 72,2 \% ), Nordeste ( 72,6 \% ) e Centro-Oeste ( 74,5 \% ) foram inferiores ao percentual nacional.
Nas demais Grandes Regiões, o percentual de alunos de Bacharelado que dispensaram Entre duas e quatro horas para concluir a prova ficou igual ou acima de 78,3 \%, como mostra o Gráfico 4.17b.
Não há diferenças estatisticamente significativas entre a proporção das Grandes Regiões.
Uma vez sendo analisadas as alternativas escolhidas pelos estudantes de Licenciatura que se situam nos diferentes quartos de desempenho, observa -se uma tendência crescente : uma maior proporção de participantes no quarto superior declarou ter gasto Entre duas e quatro horas para concluir a prova quando comparada com as dos quartos inferiores.
As diferenças entre o primeiro ( 67,9 \% ) e os demais quartos de desempenho são estatisticamente significativas.
Uma vez sendo analisadas as alternativas escolhidas pelos estudantes de Bacharelado que se situam nos diferentes quartos de desempenho, observa -se uma tendência crescente : uma maior proporção de participantes no quarto superior declarou ter gasto Entre duas e quatro horas para concluir a prova quando comparadas com as dos quartos inferiores.
As diferenças entre o quarto superior ( 82,2 \% ) e os dois quartos de desempenho inferiores, primeiro e segundo, são estatisticamente significativas.

% matched lemmas: cada, estar, lembrar, perguntar
\end{textsample}
